\section{Основы вариационного исчисления}
\label{sec:calculus-of-variations}

\subsection{Функционалы и вариационная задача}

\emph{Функционал} --- это отображение некоторого пространства функций в
числовую ось. Простейшая задача вариационного исчисления состоит в поиске
экстремума функционала
\[
    J[y]=\int_a^b f(x,y(x),y'(x))\,dx\longrightarrow \min\;(\max)
\]
на множестве допустимых функций, удовлетворяющих, например, фиксированным
краевым условиям
\[
    y(a)=y_a,\qquad y(b)=y_b.
\]
В отличие от обычной задачи на экстремум конечномерной функции здесь
неизвестным является целая функция $y(x)$.

Пусть $y(x)$ --- допустимая функция. Её вариацией называется малое изменение
\[
    y_\varepsilon(x)=y(x)+\varepsilon\eta(x),
\]
где $\varepsilon$ --- малый параметр, а $\eta(x)$ --- произвольная допустимая
вариация. При фиксированных концах
\[
    \eta(a)=\eta(b)=0.
\]
Изучение поведения $J[y+\varepsilon\eta]$ при $\varepsilon\to0$ играет ту же
роль, что исследование производных функции конечного числа переменных.

\subsection{Первая вариация}

Если функционал дифференцируем, его первая вариация определяется как линейная
по $\eta$ часть приращения:
\[
    \delta J[y;\eta]
    =\left.\frac{d}{d\varepsilon}J[y+\varepsilon\eta]
      \right|_{\varepsilon=0}.
\]
Для функционала
\[
    J[y]=\int_a^b f(x,y,y')\,dx
\]
получаем
\[
    \delta J
    =\int_a^b\left(f_y\eta+f_{y'}\eta'\right)dx.
\]
Интегрируя второй член по частям,
\[
    \delta J=
    \left.f_{y'}\eta\right|_a^b+
    \int_a^b\left(f_y-\frac{d}{dx}f_{y'}\right)\eta\,dx.
\]
Для экстремали первая вариация должна обращаться в нуль для всех допустимых
вариаций.

\subsection{Уравнение Эйлера--Лагранжа}

При фиксированных концах граничный член исчезает. Из основной леммы
вариационного исчисления следует необходимое условие экстремума:
\[
    \boxed{\quad
    f_y-\frac{d}{dx}f_{y'}=0
    \quad}
\]
--- \emph{уравнение Эйлера--Лагранжа}.

Для функционала, зависящего от нескольких функций,
\[
    J=\int_a^b
    f(x,y_1,\ldots,y_n,y_1',\ldots,y_n')\,dx,
\]
получается система
\[
    \frac{\partial f}{\partial y_i}
    -\frac{d}{dx}\frac{\partial f}{\partial y_i'}=0,
    \qquad i=1,\ldots,n.
\]
Это основное дифференциальное уравнение классического вариационного исчисления.
Его решение является лишь кандидатом на минимум или максимум: после нахождения
экстремали необходимо учитывать краевые условия и условия второго порядка.


\subsection{Естественные граничные условия и трансверсальность}

Если значение функции на одном из концов не фиксировано, соответствующая
вариация там уже не обязана обращаться в нуль. Тогда из граничного члена первой
вариации возникают дополнительные условия. Например, если $x=b$ фиксирован,
но $y(b)$ свободно, то
\[
    f_{y'}(b)=0.
\]
Такое условие называют \emph{естественным граничным условием}.

Если конец кривой может перемещаться по заданной линии, возникает более общее
условие \emph{трансверсальности}. Таким образом, вид краевых условий зависит от
того, какие данные на границе фиксированы, а какие свободны.

\subsection{Задачи с ограничениями и множители Лагранжа}

Вариационные задачи могут содержать дополнительные ограничения. Типичный пример
--- \emph{изопериметрическая задача}:
\[
    J[y]=\int_a^b f(x,y,y')\,dx\to\min,
    \qquad
    \int_a^b g(x,y,y')\,dx=K.
\]
При стандартных условиях вводится постоянный множитель Лагранжа $\lambda$ и
рассматривается расширенный функционал
\[
    \widetilde J[y]
    =\int_a^b\bigl(f+\lambda g\bigr)dx.
\]
Поэтому необходимое условие имеет вид уравнения Эйлера--Лагранжа для функции
\[
    F(x,y,y')=f(x,y,y')+\lambda g(x,y,y'):
\]
\[
    F_y-\frac{d}{dx}F_{y'}=0.
\]

При дифференциальных связях аналогично используются функции-множители
Лагранжа.

\subsection{Вторая вариация и условия минимума}

После выполнения уравнения Эйлера--Лагранжа первая вариация равна нулю, поэтому
для различения минимума и максимума рассматривают вторую вариацию. Для
простейшего функционала
\[
    \delta^2J=\int_a^b
    \left(
        f_{yy}\eta^2
        +2f_{yy'}\eta\eta'
        +f_{y'y'}(\eta')^2
    \right)dx.
\]
Необходимое условие слабого минимума:
\[
    \delta^2J\geq0
\]
для всех допустимых вариаций. Из него, в частности, следует условие Лежандра
\[
    f_{y'y'}\geq0.
\]
Одного нестрогого условия Лежандра в общем случае недостаточно. В классической
теории достаточные условия слабого минимума связывают положительную
определённость второй вариации с усиленным условием Лежандра и отсутствием
сопряжённых точек, определяемых уравнением Якоби. Для сильного минимума
дополнительно используется условие Вейерштрасса.

\subsection{Что важно сказать на экзамене}

Нужно уметь поставить задачу
$J[y]=\int f(x,y,y')dx\to\min$, ввести вариацию
$y+\varepsilon\eta$ и записать первую вариацию. После интегрирования по частям
из условия $\delta J=0$ получается уравнение Эйлера--Лагранжа
\[
    f_y-\frac{d}{dx}f_{y'}=0.
\]
При свободных концах возникают естественные условия или условия
трансверсальности. При интегральном ограничении вводится множитель Лагранжа и
уравнение Эйлера--Лагранжа записывается для $f+\lambda g$. Для проверки
минимума рассматривают вторую вариацию; простейшее необходимое условие ---
условие Лежандра $f_{y'y'}\geq0$.

\newpage
